% ----------------------------------------------------------------------
%  2.2.4. Образование дислокаций
% ----------------------------------------------------------------------
\subsection{Образование дислокаций}\label{ssec:disl-origin}

\begin{orientir}
Дислокации неравновесны (п.~\ref{ssec:disl-elastic}): тепловое движение
их не создаёт~— энергия отрезка длиной в межатомное расстояние составляет
несколько электронвольт. Откуда же они берутся? Здесь мы рассмотрим
механизмы зарождения дислокаций при кристаллизации и охлаждении, а
главное~— механизм их \emph{размножения} при пластической деформации
(источник Франка--Рида), благодаря которому плотность дислокаций при
деформации растёт на порядки.
\end{orientir}

Механизм образования дислокаций~— один из наименее изученных вопросов
теории дефектов; здесь много рабочих гипотез. Ясно одно: дислокации
возникают разными путями. Металлы сразу после кристаллизации уже содержат
много дислокаций~— значит, они рождаются у фронта кристаллизации или при
последующем охлаждении.

\subsubsection{Дислокации при росте кристалла}

\paragraph{Винтовая дислокация облегчает рост}
Совершенный кристалл растёт, присоединяя атомы к ступеньке на
поверхности. Но когда слой достроен, ступенька исчезает
(рис.~\ref{fig:growth},\,а), и для нового слоя нужно зародить на гладкой
грани двумерный «островок»~— это самое трудное звено роста, оно требует
больших пересыщений (переохлаждений). Если же на грань выходит винтовая
дислокация, ступенька никогда не исчезает (рис.~\ref{fig:growth},\,б):
атомы укладываются на винтовую поверхность, ступенька вращается вокруг
точки выхода дислокации, и рост идёт непрерывно, по спирали.

\begin{figure}[ht]
\centering
\begin{subfigure}[b]{0.42\textwidth}\centering
\begin{tikzpicture}[cubeproj, scale=0.9, line join=round]
  \filldraw[fill=fmBlue!20, draw=atomLine] (3,0,0)--(3,2.5,0)--(3,2.5,1.6)--(3,0,1.6)--cycle;
  \filldraw[fill=fmBlue!8, draw=atomLine] (0,0,1.6)--(3,0,1.6)--(3,2.5,1.6)--(0,2.5,1.6)--cycle;
  \filldraw[fill=fmBlue!14, draw=atomLine] (0,0,0)--(3,0,0)--(3,0,1.6)--(0,0,1.6)--cycle;
  \filldraw[fill=fmAmber!30, draw=atomLine] (0,0,1.6)--(1.4,0,1.6)--(1.4,0,1.9)--(0,0,1.9)--cycle;
  \filldraw[fill=fmAmber!15, draw=atomLine] (0,0,1.9)--(1.4,0,1.9)--(1.4,2.5,1.9)--(0,2.5,1.9)--cycle;
  \filldraw[fill=fmAmber!45, draw=atomLine] (1.4,0,1.6)--(1.4,2.5,1.6)--(1.4,2.5,1.9)--(1.4,0,1.9)--cycle;
  \node[atomB=2.6mm] at (1.6,1.2,1.75) {};
  \draw[arr, fmRed] (2.6,1.2,2.6) -- (1.75,1.2,1.95);
\end{tikzpicture}
\caption{}
\end{subfigure}\hfill
\begin{subfigure}[b]{0.42\textwidth}\centering
\begin{tikzpicture}[scale=0.75]
  \fill[fmBlue!6] (-2.6,-2.6) rectangle (2.6,2.6);
  \draw[fmAmber!90!black, line width=1.2pt, domain=0:4.6*pi, samples=300, smooth, variable=\t]
     plot ({0.17*\t*cos(\t r)}, {0.17*\t*sin(\t r)});
  \fill[fmBlue] (0,0) circle (2.2pt);
  \node[slbl, fmBlue, below right] at (0,0) {выход дислокации};
\end{tikzpicture}
\caption{}
\end{subfigure}
\caption{Рост кристалла: \textit{а}~— ступенька на поверхности
совершенного кристалла исчезает, когда слой достроен; \textit{б}~— вид
сверху на грань, на которую выходит винтовая дислокация: ступенька
закручивается в спираль и никогда не исчезает}
\label{fig:growth}
\end{figure}

Но как возникает такое закручивание в самом начале роста? Кристаллы,
как правило, зарождаются не самопроизвольно, а \emph{гетерогенно}~— на
стенках изложницы или на мельчайших твёрдых частицах в расплаве. На
поверхность таких подложек выходят винтовые дислокации; к готовым
ступенькам присоединяются атомы, и дислокация из подложки «прорастает»
в кристалл.

\paragraph{Дислокации несоответствия}
При ориентированном наращивании кристалла на подложку (\term{эпитаксии})
параметры двух решёток всегда немного различаются. Пока растущий слой
тонкий, несоответствие компенсируется упругой деформацией. Когда толщина
слоя превышает критическую, выгоднее снять часть напряжений, введя на
границе дислокации (рис.~\ref{fig:misfit}). Их называют
\term{дислокациями несоответствия} (структурными, эпитаксиальными). Чем
больше несоответствие решёток, тем чаще расположены такие дислокации:
энергия, затраченная на их образование, окупается уменьшением энергии
упругой деформации.

\begin{figure}[ht]
\centering
\begin{tikzpicture}[scale=0.5]
  \foreach \i in {0,...,11}{\foreach \j in {0,...,2}{\node[atom=3.6mm] at (\i,\j) {};}}
  \foreach \i in {0,...,12}{\foreach \j in {3,...,5}{
    \pgfmathsetmacro{\xx}{\i*11/12}
    \node[atom=3.6mm, fill=fmGreen!40] at (\xx,\j+0.0) {};}}
  \dislT[fmRed]{(5.5,2.45)}{0}
  \node[lbl, right] at (12.4,4) {кристалл К};
  \node[lbl, right] at (12.4,1) {подложка П};
  \draw[fmGray, dashed] (-0.6,2.5) -- (12,2.5);
\end{tikzpicture}
\caption{Дислокация несоответствия на границе растущего кристалла К с
подложкой П: на 12 межатомных расстояний подложки приходится 13
расстояний кристалла, и одна атомная плоскость кристалла оканчивается на
границе (схема без релаксации)}
\label{fig:misfit}
\end{figure}

Похожая ситуация возникает из-за \emph{ликвации} (сегрегации примесей) при
кристаллизации: соседние слои имеют разный состав и немного разные
межатомные расстояния, и их сопряжение может пойти через структурные
дислокации.

\paragraph{Мозаичная структура}
Случайности роста (например, изгиб тонких ветвей дендрита конвекционными
потоками и градиентами температуры) приводят к тому, что соседние части
кристалла слегка разориентированы. Когда они срастаются, на границе
возникают дислокации. Простейший случай~— срастание двух симметрично
разориентированных частей: на границе образуется стенка краевых
дислокаций (п.~\ref{ssec:disl-structures}, рис.~\ref{fig:tilt-wall}).
Кристалл получается состоящим из слегка разориентированных
\term{субзёрен} (блоков)~— это \term{мозаичная структура}.

\subsubsection{Дислокации в затвердевшем металле}

\paragraph{Захлопывание вакансионных дисков}
Основным механизмом считают вакансионный. Равновесная концентрация
вакансий резко падает при охлаждении от температуры кристаллизации, и при
быстром охлаждении кристалл сильно пересыщается вакансиями. Избыточные
вакансии объединяются в плоские диски, параллельные плоскостям плотнейшей
упаковки (толщиной в один-три слоя вакансий). Когда диаметр диска
превышает критический, его стороны под действием межатомного притяжения
сближаются, и диск \emph{захлопывается}, образуя дислокационную петлю
(рис.~\ref{fig:vac-disk}). То же происходит при охлаждении после нагрева
под закалку, а при облучении петли возникают на плоских скоплениях и
вакансий, и межузельных атомов.

\begin{figure}[ht]
\centering
\begin{subfigure}[b]{0.46\textwidth}\centering
\begin{tikzpicture}[scale=0.45]
  \foreach \i in {-6,...,6}{\foreach \j in {-2,...,2}{
    \pgfmathtruncatemacro{\hole}{(\j==0 && \i>-3 && \i<3) ? 1 : 0}
    \ifnum\hole=1 \node[vac=3.4mm] at (\i,\j) {};\else\node[atom=3.4mm] at (\i,\j) {};\fi}}
\end{tikzpicture}
\caption{}
\end{subfigure}\hfill
\begin{subfigure}[b]{0.46\textwidth}\centering
\begin{tikzpicture}[scale=0.45]
  \foreach \i in {-6,...,6}{\foreach \j in {-2,...,2}{
    \pgfmathsetmacro{\dip}{0.5*(1/(1+exp(min(9,-(\i+2.5)*1.6))))*(1/(1+exp(min(9,(\i-2.5)*1.6))))}
    \pgfmathsetmacro{\yy}{(\j>0 ? \j-\dip*(1-0.25*(\j-1)) : (\j<0 ? \j+\dip*(1-0.25*(-\j-1)) : \j))}
    \pgfmathtruncatemacro{\hole}{(\j==0 && \i>-3 && \i<3) ? 1 : 0}
    \ifnum\hole=0 \node[atom=3.4mm] at (\i,\yy) {};\fi}}
  \dislT[fmRed]{(-2.45,0)}{90}
  \dislT[fmRed]{(2.45,0)}{-90}
\end{tikzpicture}
\caption{}
\end{subfigure}
\caption{Образование дислокационной петли из диска вакансий (сечение,
перпендикулярное диску): \textit{а}~— диск вакансий в атомной плоскости;
\textit{б}~— диск захлопнулся, соседние плоскости сомкнулись; по краю
бывшего диска проходит петля краевой дислокации (в сечении~— две
дислокации противоположного знака)}
\label{fig:vac-disk}
\end{figure}

\paragraph{Концентрация напряжений}
Дислокации зарождаются и там, где напряжения локально достигают
$\sim G/30$: около включений, трещин, границ двойников. Например, при
охлаждении включение и матрица сжимаются по-разному, и у их границы
возникают напряжения, достаточные для самопроизвольного зарождения
дислокационных петель. Петли уходят от включения, и напряжения
релаксируют.

\subsubsection{Размножение дислокаций: источник Франка--Рида}

При холодной пластической деформации плотность дислокаций растёт на
несколько порядков (от $\sim10^{6}$--$10^{8}$ до $10^{11}$--$10^{12}$~см$^{-2}$).
Значит, должен существовать механизм \emph{размножения} дислокаций. Его
предложили Франк и Рид (1950).

\paragraph{Линейное натяжение}
Энергия дислокации пропорциональна её длине, поэтому дислокация
стремится сократиться, как натянутая струна. Это \term{линейное
натяжение}~— приращение энергии на единицу длины линии:
\begin{equation}\label{eq:line-tension}
T\approx\alpha Gb^2,\qquad\alpha\approx\tfrac12 .
\end{equation}
Чтобы дислокация оставалась изогнутой, нужна сила, которая уравновешивает
стремление выпрямиться. Её создаёт касательное напряжение
(рис.~\ref{fig:bowed}).

\begin{figure}[ht]
\centering
\begin{tikzpicture}[scale=1.0]
  \coordinate (O) at (0,0);
  \draw[disl] (60:3) arc (60:120:3);
  \draw[fmGray, dashed] (O) -- (60:3) (O) -- (120:3) (O) -- (90:3);
  \draw (80:0.9) arc (80:100:0.9); \node[slbl] at (90:1.15) {$\d\theta$};
  \node[slbl, below] at (O) {$O$};
  \draw[arr, fmBlue] (120:3) -- ++(210:0.9) node[left, slbl]{$T$};
  \draw[arr, fmBlue] (60:3) -- ++(-30:0.9) node[right, slbl]{$T$};
  \draw[arr, fmRed, line width=1pt] (90:3) -- (90:3.9) node[above, slbl]{$\tau b\,\d S$};
  \node[slbl, right] at (75:1.8) {$r$};
  \node[slbl, above right] at (75:3.05) {$\d S$};
\end{tikzpicture}
\caption{Равновесие изогнутого элемента дислокации: силы линейного
натяжения $T$ на его концах дают равнодействующую $T\,\d\theta$,
направленную к центру кривизны; её уравновешивает сила $\tau b\,\d S$}
\label{fig:bowed}
\end{figure}

\begin{shagi}
\item Элемент дуги длиной $\d S$ с радиусом кривизны $r$ стягивает угол
$\d\theta=\d S/r$.
\item Силы натяжения $T$ на его концах направлены по касательным; их
равнодействующая $2T\sin(\d\theta/2)\approx T\,\d\theta$ направлена к
центру кривизны.
\item Касательное напряжение действует на элемент с силой $\tau b\,\d S$
наружу.
\item Равновесие: $T\,\d\theta=\tau b\,\d S$, откуда $T=\tau br$.
\end{shagi}
\begin{keyeq}\label{eq:tau-r}
\tau=\frac{T}{br}=\frac{\alpha Gb}{r} .
\end{keyeq}
Это напряжение, необходимое, чтобы изогнуть дислокацию до радиуса $r$.

\paragraph{Как работает источник}
Пусть отрезок дислокации $AB$ длиной $l$ закреплён на концах (например,
в узлах дислокационной сетки или на частицах) и лежит в своей плоскости
скольжения (рис.~\ref{fig:frank-read}).

\begin{figure}[ht]
\centering
\begin{tikzpicture}[scale=0.62]
  \tikzset{pin/.style={circle, fill=black, inner sep=1.6pt}, fr/.style={disl, line width=1.3pt}}
  \foreach \s/\lab in {0/а,5.5/б,11/в}{
    \begin{scope}[shift={(\s,0)}]
    \node[lbl] at (0,-1.6) {\lab};
    \end{scope}}
  % а
  \begin{scope}[shift={(0,0)}]
    \fill[fmAmber!0] (-1,0) -- (1,0);
    \draw[fr] (-1,0) -- (1,0);
    \node[pin] at (-1,0) {}; \node[pin] at (1,0) {};
    \node[slbl, below] at (-1,-0.1) {$A$}; \node[slbl, below] at (1,-0.1) {$B$};
    \draw[arr, fmRed] (0,0.2) -- (0,0.9) node[above, slbl]{$\tau$};
  \end{scope}
  % б
  \begin{scope}[shift={(5.5,0)}]
    \fill[fmAmber!30] (-1,0) arc (180:0:1) -- cycle;
    \draw[fr] (-1,0) arc (180:0:1);
    \node[pin] at (-1,0) {}; \node[pin] at (1,0) {};
    \node[slbl, above] at (0,1.0) {$r=l/2$};
  \end{scope}
  % в
  \begin{scope}[shift={(11,0)}]
    \fill[fmAmber!30] (1,0) arc (-31:211:1.166) -- cycle;
    \draw[fr] (1,0) arc (-31:211:1.166);
    \node[pin] at (-1,0) {}; \node[pin] at (1,0) {};
  \end{scope}
  % г
  \begin{scope}[shift={(0,-5.2)}]
    \fill[fmAmber!30] plot[smooth] coordinates {(-1,0) (-0.55,-0.75) (-0.12,-1.05) (-0.9,-1.5) (-2.0,-0.8) (-2.2,0.6) (-1.2,2.0) (0,2.3) (1.2,2.0) (2.2,0.6) (2.0,-0.8) (0.9,-1.5) (0.12,-1.05) (0.55,-0.75) (1,0)} -- cycle;
    \draw[fr] plot[smooth] coordinates {(-1,0) (-0.55,-0.75) (-0.12,-1.05) (-0.9,-1.5) (-2.0,-0.8) (-2.2,0.6) (-1.2,2.0) (0,2.3) (1.2,2.0) (2.2,0.6) (2.0,-0.8) (0.9,-1.5) (0.12,-1.05) (0.55,-0.75) (1,0)};
    \node[pin] at (-1,0) {}; \node[pin] at (1,0) {};
    \draw[fmRed, thick] (0,-1.05) circle (0.32);
    \node[lbl] at (0,-2.6) {г};
  \end{scope}
  % д
  \begin{scope}[shift={(5.5,-5.2)}]
    \fill[fmAmber!30] plot[smooth cycle] coordinates {(0,-1.5) (-1.4,-1.35) (-2.2,-0.4) (-2.25,0.8) (-1.2,2.05) (0,2.35) (1.2,2.05) (2.25,0.8) (2.2,-0.4) (1.4,-1.35)};
    \fill[white] (-1,0) .. controls (-0.5,-0.55) and (0.5,-0.55) .. (1,0) -- cycle;
    \draw[fr] plot[smooth cycle] coordinates {(0,-1.5) (-1.4,-1.35) (-2.2,-0.4) (-2.25,0.8) (-1.2,2.05) (0,2.35) (1.2,2.05) (2.25,0.8) (2.2,-0.4) (1.4,-1.35)};
    \draw[fr] (-1,0) .. controls (-0.5,-0.55) and (0.5,-0.55) .. (1,0);
    \node[pin] at (-1,0) {}; \node[pin] at (1,0) {};
    \node[lbl] at (0,-2.6) {д};
  \end{scope}
  % е
  \begin{scope}[shift={(11,-5.2)}]
    \draw[fr] plot[smooth cycle] coordinates {(0,-1.9) (-1.8,-1.6) (-2.6,-0.4) (-2.6,0.9) (-1.4,2.4) (0,2.7) (1.4,2.4) (2.6,0.9) (2.6,-0.4) (1.8,-1.6)};
    \draw[fr] (-1,0) -- (1,0);
    \node[pin] at (-1,0) {}; \node[pin] at (1,0) {};
    \draw[arr, fmGray] (0,2.75) -- (0,3.3);
    \draw[arr, fmGray] (2.65,0.3) -- (3.2,0.3);
    \draw[arr, fmGray] (-2.65,0.3) -- (-3.2,0.3);
    \node[lbl] at (0,-2.6) {е};
  \end{scope}
\end{tikzpicture}
\caption{Работа источника Франка--Рида (вид на плоскость скольжения;
закрашена площадь, заметённая дислокацией): \textit{а}~— исходный отрезок
$AB$; \textit{б}~— дуга выгнулась в полуокружность (критическое
состояние); \textit{в}~— петля неустойчиво расширяется;
\textit{г}~— петля обвивает точки закрепления, и её участки
противоположного знака сближаются (кружок); \textit{д}~— они
аннигилировали: образовалась замкнутая петля и новый отрезок $AB$;
\textit{е}~— петля уходит, отрезок выпрямляется, и цикл повторяется}
\label{fig:frank-read}
\end{figure}

\begin{enumerate}[label=\arabic*)]
\item Под действием напряжения $\tau$ отрезок выгибается в дугу. По
формуле~\eqref{eq:tau-r} каждому $\tau$ соответствует свой радиус
кривизны; пока $\tau$ мало, дуга устойчива, и при снятии нагрузки линейное
натяжение упруго её выпрямляет.
\item С ростом $\tau$ радиус уменьшается. Наименьший возможный радиус~—
у полуокружности: $r=l/2$. Ему отвечает наибольшее
напряжение\ispr{В оригинале (с.~260) формула записана как
$\tau_{\text{кр}}=2\alpha GbL/l$; лишний множитель $L$~— опечатка.}:
\begin{keyeq}\label{eq:tau-FR}
\tau_{\text{кр}}=\frac{2\alpha Gb}{l}\approx\frac{Gb}{l}\qquad(\alpha=\tfrac12) .
\end{keyeq}
\item При любом небольшом превышении $\tau_{\text{кр}}$ дальнейшее
расширение петли \emph{увеличивает} радиус кривизны, а значит, требует
меньшего напряжения. Положение становится неустойчивым, и петля
расширяется сама, заметая всё бо́льшую площадь.
\item Петля закреплена в $A$ и $B$ и потому «закручивается» вокруг этих
точек. Наконец два её участка за отрезком $AB$ встречаются. Здесь сходятся
винтовые участки \emph{противоположного знака}~— они аннигилируют.
\item Дислокация распадается на замкнутую петлю и новый отрезок $AB$.
Петля свободно расширяется дальше, а отрезок под действием напряжения и
натяжения возвращается в исходное положение~— и всё повторяется.
\end{enumerate}

\begin{glavnoe}
Источник Франка--Рида может испускать неограниченное число петель в
одной плоскости скольжения и создавать в ней большой сдвиг. Для его
работы нужно напряжение $\tau_{\text{кр}}\approx Gb/l$: чем длиннее
закреплённый отрезок, тем легче работает источник.
\end{glavnoe}

\begin{smysl}[Почему дислокация так странно изгибается]
Напряжение всё время действует в одном направлении, а петля растёт во все
стороны. Дело в том, что изгибаясь, исходно краевой отрезок становится
смешанным: на нём появляются краевые, винтовые и смешанные участки.
Сила $\tau b$ всегда перпендикулярна линии (п.~\ref{ssec:disl-elastic}),
и каждый участок движется по нормали к себе: краевые~— вдоль $\vb b$,
винтовые~— поперёк. В итоге петля расширяется во все стороны.
\end{smysl}

\paragraph{Пределы работы источника}
Петли, испущенные источником, не уходят за пределы зерна: они
останавливаются у границы и образуют \term{скопление}
(п.~\ref{ssec:disl-structures}). Поле скопления действует на источник
обратным напряжением, и когда оно сравнивается с приложенным, источник
останавливается. У границ зёрен наблюдали серии из нескольких десятков
дислокаций от одного источника. Источник может и «развалиться», испустив
10--20 дислокаций: поле приводит в движение отрезок, служивший ему
опорой.

Если вблизи источника есть непреодолимое препятствие $C$
(рис.~\ref{fig:FR-variants},\,а), критическое напряжение определяется не
длиной $AB$, а меньшим расстоянием до препятствия. Поэтому длина базы
источника внутри зерна не превышает $\sim1/3$ размера зерна $d$.
\term{Однополюсный} источник (дислокация закреплена в одной точке $A$, а
другим концом выходит на поверхность; рис.~\ref{fig:FR-variants},\,б)
вращается вокруг точки закрепления, наматываясь в спираль. Его база может
достигать $d/2$, а критическое напряжение соответственно меньше.

\begin{figure}[ht]
\centering
\begin{subfigure}[b]{0.42\textwidth}\centering
\begin{tikzpicture}[scale=0.7]
  \draw[fmGray, line width=1pt] (-2.5,-1.2) -- (-2.5,2.4) node[above, slbl]{граница};
  \draw[disl] (-1,0) -- (1.5,0);
  \draw[disl, dashed] (-1,0) arc (180:0:1.25);
  \node[circle, fill=black, inner sep=1.5pt] at (-1,0) {}; \node[circle, fill=black, inner sep=1.5pt] at (1.5,0) {};
  \node[slbl, below] at (-1,0) {$A$}; \node[slbl, below] at (1.5,0) {$B$};
  \fill[black] (-1.8,0.9) circle (2.2pt) node[left, slbl]{$C$};
\end{tikzpicture}
\caption{}
\end{subfigure}\hfill
\begin{subfigure}[b]{0.42\textwidth}\centering
\begin{tikzpicture}[scale=0.7]
  \fill[fmBlue!8] (0,-2) rectangle (2.6,2);
  \draw[fmGray, line width=1pt] (2.6,-2) -- (2.6,2) node[above, slbl]{поверхность};
  \draw[disl, domain=0:3.2*pi, samples=200, smooth, variable=\t] plot ({0.07*\t*cos(\t r)}, {0.07*\t*sin(\t r)});
  \draw[disl] (-0.15,-0.68) .. controls (1.2,-1.0) .. (2.6,-0.6);
  \node[circle, fill=black, inner sep=1.5pt] at (0,0) {};
  \node[slbl, left] at (-0.1,0.1) {$A$};
\end{tikzpicture}
\caption{}
\end{subfigure}
\caption{Варианты источников: \textit{а}~— двухполюсный источник вблизи
препятствия $C$: критическое напряжение определяется расстоянием до $C$;
\textit{б}~— однополюсный источник: дислокация закреплена в точке $A$,
другим концом выходит на поверхность и наматывается в спираль}
\label{fig:FR-variants}
\end{figure}

\paragraph{Размножение двойным поперечным скольжением}
Источник Франка--Рида может возникнуть и при двойном поперечном
скольжении (рис.~\ref{fig:double-cross}). Винтовой участок дислокации,
остановленный препятствием, переходит в плоскость поперечного скольжения,
а затем в параллельную исходной плоскость. Если касательное напряжение в
плоскости поперечного скольжения мало, оставшиеся там участки
малоподвижны и служат точками закрепления: отрезок в новой плоскости
становится источником Франка--Рида. \emph{Множественное} поперечное
скольжение даёт длинную дислокационную линию, переходящую из плоскости в
плоскость, с источниками в каждой из них,~— очень эффективный механизм
размножения, объясняющий образование широких полос скольжения.

\begin{figure}[ht]
\centering
\begin{tikzpicture}[cubeproj, scale=1.0, line join=round]
  \def\Lx{6}
  \filldraw[fill=fmBlue!7, draw=fmBlue!45] (0,0,0) -- (\Lx,0,0) -- (\Lx,2,0) -- (0,2,0) -- cycle;
  \filldraw[fill=fmGreen!10, draw=fmGreen!55] (0,2,0) -- (\Lx,2,0) -- (\Lx,2.9,1.4) -- (0,2.9,1.4) -- cycle;
  \filldraw[fill=fmBlue!7, draw=fmBlue!45] (0,2.9,1.4) -- (\Lx,2.9,1.4) -- (\Lx,5.0,1.4) -- (0,5.0,1.4) -- cycle;
  \fill[fmAmber!35] (1.9,2.9,1.4) .. controls (1.7,4.9,1.4) and (4.3,4.9,1.4) .. (4.1,2.9,1.4) -- cycle;
  \draw[disl, line width=1.6pt] (0,1.5,0) -- (1.6,1.5,0) -- (1.75,2.45,0.7) -- (1.9,2.9,1.4)
       .. controls (1.7,4.9,1.4) and (4.3,4.9,1.4) .. (4.1,2.9,1.4) -- (4.25,2.45,0.7) -- (4.4,1.5,0) -- (\Lx,1.5,0);
  \fill (1.9,2.9,1.4) circle (2pt); \fill (4.1,2.9,1.4) circle (2pt);
  \draw[bvec] (0.4,0.4,0) -- (1.6,0.4,0) node[below right, slbl]{$\vb b$};
  \node[slbl, fmBlue] at (5.4,0.3,0) {$P$};
  \node[slbl, fmGreen] at (5.6,2.5,0.75) {$R$};
  \node[slbl, fmBlue] at (5.4,4.7,1.4) {$S$};
\end{tikzpicture}
\caption{Источник, возникший при двойном поперечном скольжении: участки
дислокации в плоскости поперечного скольжения $R$ малоподвижны и служат
точками закрепления; отрезок между ними выгибается в плоскости $S$, как
источник Франка--Рида}
\label{fig:double-cross}
\end{figure}

\begin{itogi}
\begin{itemize}
\item Дислокации не рождаются тепловыми флуктуациями. Они возникают при
росте кристаллов (винтовые дислокации из подложки, дислокации
несоответствия, срастание разориентированных частей), при захлопывании
вакансионных дисков, у концентраторов напряжений ($\sim G/30$).
\item Линейное натяжение $T\approx\frac12Gb^2$; изогнуть дислокацию до
радиуса $r$ стоит напряжения $\tau=\alpha Gb/r$.
\item Источник Франка--Рида: закреплённый отрезок длиной $l$ начинает
непрерывно испускать петли при $\tau_{\text{кр}}\approx Gb/l$.
\item Размножение идёт также через однополюсные источники и двойное
поперечное скольжение.
\end{itemize}
\end{itogi}

\begin{proverka}
\item Почему кристалл с винтовой дислокацией растёт при меньших
пересыщениях, чем совершенный?
\item Почему при большом несоответствии решёток дислокации несоответствия
расположены чаще?
\item Почему источник Франка--Рида «срабатывает» именно при радиусе
кривизны $l/2$?
\item Оцените $\tau_{\text{кр}}$ для источника длиной $l=1$~мкм в меди
($G=48$~ГПа, $b=0{,}26$~нм).
\item Почему источник прекращает работу, даже если нагрузка сохраняется?
\end{proverka}
